%----------------------------------------------------------------
% Abstract
%----------------------------------------------------------------
\chapter*{Abstract}
\addcontentsline{toc}{chapter}{Abstract}

Water distribution networks deliver drinking water to millions of
people daily, and the sensor arrays attached to their
Supervisory Control and Data Acquisition (SCADA) systems are the only real-time window operators
have into network state. That same connectivity is what makes them
worth attacking: an adversary who reaches even a few pressure or
flow sensors can inject readings that look, to any classical
statistical monitor, indistinguishable from measurement noise.

This thesis builds a family of Graph Neural Network (GNN) models that jointly
fill in missing sensor readings and flag compromised sensors. The
work advances in four stages. First,
a spatial multi-task \gnn{} baseline shows that graph structure alone
encodes enough hydraulic context to beat classical reconstruction
methods by more than an order of magnitude in MAE. Second,
a temporal extension wraps each expert in a Gated Recurrent Unit (GRU) over a
six-step sliding window, which is the minimum history needed to flag
replay attacks, where an attacker broadcasts a legitimate past reading
to hide a real anomaly. Third, a Mixture of Experts (MoE) adds a learned router
that assigns each snapshot to the specialist most likely to recognise
its attack type, improving aggregate anomaly F1 by around four
percentage points on the Modena benchmark (272 nodes, 317 pipes).

The fourth and principal contribution is an adversarial self-play
loop. An \emph{Attacker} \gnn{} co-evolves against the defender in a
Stackelberg game: it learns to produce sparse, budget-respecting,
physics-aware perturbations that fool the current defender, while the
defender is updated against those perturbations to catch them.
Extending the attacker to a Mixture of Experts (MoE) of specialists lets the
population auto-discover two functional attack families, labelled
\emph{bold} and \emph{subtle}, without any attack-class supervision.
The resulting self-play defender achieves an anomaly F1 of 0.767 on
Modena, a gain of 5.8 percentage points over the pretrained baseline,
and reduces pressure reconstruction error by 23\%. Most importantly,
the self-play defender outperforms the pretrained model on two attack
types held out from training entirely, demonstrating genuine
adversarial robustness rather than curve-fitting to a fixed attack
distribution.

Experiments cover three EPANET benchmark networks spanning
11 to 272 junctions, five hand-crafted attack types, two novel
held-out attacks, a three-seed reproducibility study, and comparisons
against classical state-estimation baselines. An interactive
Streamlit dashboard accompanies the thesis and makes the models,
results, and explainability analyses accessible without running
training code.

\vspace{1em}
\noindent\textbf{Keywords:} water distribution network, graph neural
network, anomaly detection, cyber-physical attack, adversarial self-play,
mixture of experts, Stackelberg game, temporal modeling.
